% status: FULL
\chapter{Causal Attention, One Row at a Time}\label{ch:causal-attention}
\begin{ChapterIntent}
A query, a key and a value are useful only when we connect them. This chapter
does that with two tokens and two-dimensional vectors. We will calculate
one attention output by hand, prevent it from reading the future, and turn
the calculation into a small program. Only after the result is clear will
we write the compact matrix equation used in papers.
\end{ChapterIntent}

\section{The missing operation is a weighted sum}
Imagine that two earlier positions contain information useful to the current
token. One carries a value vector $[2,0]$; the other carries $[0,4]$. These
are illustrative numbers, not representations extracted from a trained
model. If the current position gives the first value one quarter of its
attention and the second three quarters, its output is
\[
 \tfrac14[2,0]+\tfrac34[0,4]=[0.5,3].
\]
Notice what happened. We did not choose one value and discard the other.
We formed a new vector from both. Its first coordinate came from the first
position and its second mostly from the second position. Attention is a
learned way to produce the weights for this mixture.

The query says what the current position is looking for. Each key supplies
the information against which that query is compared. Each value supplies
the information that can be carried forward. These descriptions are useful
roles, not a claim that a vector coordinate has a human-readable meaning.
The previous chapters constructed Q, K and V using learned projections.
This chapter consumes their numerical outputs.

For now, use one head and one sequence. Each query and key has two
coordinates, so the head width is $d_h=2$. Values also have two coordinates
in this example. A real architecture may choose another value width, but
every value in this one mixture must have the same width.

\section{Turn comparisons into probabilities}
Take the positioned query and keys
\[
 \mathbf q=[1,0],\quad
 \mathbf k_0=[0,0],\quad
 \mathbf k_1=[\sqrt2\log3,0].
\]
The dot product multiplies corresponding coordinates and adds them.
Attention divides that result by $\sqrt{d_h}$. Our two scores are therefore
\[
 s_0=\frac{1\cdot0+0\cdot0}{\sqrt2}=0,\qquad
 s_1=\frac{1\cdot\sqrt2\log3+0\cdot0}{\sqrt2}=\log3.
\]
Scores are not yet weights: they may be negative and need not sum to one.
Softmax exponentiates and normalizes them. Exponentiating $[0,\log3]$
gives $[1,3]$. Dividing by their sum, four, gives $[1/4,3/4]$.
We have recovered the mixture at the start of the chapter.

Why divide by the square root of the head width? Under an illustrative
assumption of independent, zero-mean query and key coordinates with unit
variance, the dot product has variance $d_h$. Scaling by $1/\sqrt{d_h}$
keeps its variance near one. Learned activations need not satisfy those
assumptions exactly. The argument explains the design motivation; the model
definition specifies the actual scale. Omitting it silently changes the
model's computation.

\EngineFigure{foundation-attention-row}{An illustrative attention row. The
query produces scores, softmax converts scores to mixing weights, and those
weights combine values. The arrows carry the displayed numerical objects;
they do not represent a hard search or a retrieved text passage.}
{fig:foundation-attention-row}

The compact formula now has a concrete interpretation:
\[
 a_j=\frac{\exp(s_j)}{\sum_t\exp(s_t)},\qquad
 \mathbf y=\sum_j a_j\mathbf v_j.
\]
The index $j$ ranges over permitted key positions. A query produces one
row of weights, then one output vector. It does not produce a probability
distribution over vocabulary tokens. That later distribution is produced
by the model's final output projection.

\section{Do not let a token read its future}
When predicting text, position zero must not use a value from position one.
The later token would reveal information not available when the earlier
prediction was made. A \emph{causal mask} enforces this rule.

Suppose both example tokens are in a known prompt. We can compute their
projections together, but the first query may read only the first key and
value. Its permitted score row has one entry. Softmax of a singleton is
$[1]$, regardless of the finite score, so its attention output is exactly
$[2,0]$. The second query can read both positions and produces $[0.5,3]$.
Known input is not permission to leak future information into earlier rows.

There are two common ways to express the mask. A dense implementation sets
forbidden scores to negative infinity before softmax; their exponentials
contribute zero. Our small program avoids constructing forbidden scores at
all: query $i$ receives only keys and values through position $i$.
These approaches represent the same permitted set.

Masking \emph{after} normalization is wrong unless you renormalize the
remaining probabilities. In our two-entry example, zeroing the second
weight after softmax leaves $[1/4,0]$. The first query would output
$[0.5,0]$, not $[2,0]$. The output is finite, has the right shape, and is
still incorrect. This is why a shape test alone cannot validate attention.

Every row in ordinary causal self-attention can attend to itself, so its
permitted set is nonempty. Padding, custom masks or a programming error can
produce a row with no permitted keys. There is no normalized distribution
over an empty set. Our teaching routine rejects that input rather than
pretending the output is an ordinary attention result.

\section{The diagonal is permitted because the target is shifted}
A common first question is whether attending to the current token is itself
cheating. Write an illustrative input sequence as $[a,b,c]$. The logit row
at position zero predicts what comes \emph{after} $a$, the row at position
one predicts what comes after $[a,b]$, and the row at position two predicts
what comes after $[a,b,c]$. In a next-token training example their targets
would be $[b,c,d]$, where $d$ is the following token. Input position $i$
contains a known token; its target is the next token. Allowing $j=i$
therefore does not reveal that target.

This is the distinction between a token's position and a prediction's
boundary. During inference, the final prompt token is known, so its key
and value belong in the history used to predict the first continuation.
If we remove the diagonal, the first query has no permitted key. If we
remove the future mask instead, earlier queries can read their own targets
from later input rows. Neither change implements the stated decoder.

The same reasoning explains prompt parallelism. At one layer, all input
rows to that layer already exist. Their projections can be calculated
together. Each attention row then reads a different permitted subset.
Parallel calculation does not mean unrestricted information flow.
The next layer starts after the current layer has produced its outputs;
there remains a dependency across layers even when rows within a layer
are processed together.

\section{Evaluate softmax without overflowing}
An exponential grows quickly. Directly computing $\exp(1000)$ can overflow
even when the desired probability is perfectly ordinary. Subtract the
largest permitted score $m$ before exponentiating:
\[
 a_j=\frac{\exp(s_j-m)}{\sum_t\exp(s_t-m)}.
\]
The common factor cancels in real arithmetic. At least one exponent is
zero, and the others are nonpositive. For $[1000,999]$, the routine
evaluates $[1,e^{-1}]$ and normalizes it, giving approximately
$[0.731059,0.268941]$. The computation becomes numerically practical
without changing the intended distribution.

This fixes one numerical problem, not every problem. Nonfinite inputs,
rounding in the score reduction and underflow of very small probabilities
still need a contract. Our routine accepts a nonempty row of finite,
already-permitted scores. Masking belongs to its caller. This narrow
interface makes its tests easy to understand.

\section{The matrix equation is many copies of this row}
For $S$ positions, put queries, keys and values in matrices
$\mathbf Q,\mathbf K,\mathbf V\in\mathbb R^{S\times d_h}$. Then
$\mathbf Q\mathbf K^{\mathsf T}$ contains every query--key dot product.
Entry $(i,j)$ compares query $i$ with key $j$. Apply the causal mask
$M_{ij}=0$ for $j\le i$ and $M_{ij}=-\infty$ otherwise:
\[
 \mathbf A=\operatorname{softmax}_{\rm row}
 \left(\frac{\mathbf Q\mathbf K^{\mathsf T}}{\sqrt{d_h}}+\mathbf M\right),
 \qquad \mathbf Y=\mathbf A\mathbf V.
\]
Softmax acts across each row, not down each column. The output has shape
$S\times d_h$. This is the scaled dot-product construction associated with
the Transformer \cite{vaswani2017attention}; the causal choice of permitted
keys makes it suitable for this decoder.

Several heads repeat the operation on different projected coordinates.
With grouped-query attention, more than one query head uses the same key
and value head. For $H_q=4$ and $H_{kv}=2$, consecutive query heads map to
KV heads $[0,0,1,1]$. Sharing keys does not make the queries equal, so their
attention rows can still differ. Concatenating head outputs and applying an
output projection returns to model width; it does not add the heads' query
scores together.

\section{A cached chunk needs positions, not a square triangle}
Now keep three previously evaluated positions and process two additional
tokens together. There are two query rows but five key rows. A causal
mask with two rows and five columns is not the upper-left corner of the
ordinary five-by-five triangle. Query row zero belongs to logical position
three, not position zero. It may read keys zero through three. Query row
one belongs to logical position four and may read all five keys.

Choose both queries to be $[0,0]$ and every key to be $[1,2]$. All scores
are zero regardless of scaling. Give the five values one coordinate each:
$[0],[1],[2],[3],[4]$. The correct outputs are the means
\[
 y_3=(0+1+2+3)/4=1.5,\qquad
 y_4=(0+1+2+3+4)/5=2.
\]
An upper-left triangle would admit only the first key for row zero and
the first two for row one. It returns $0$ and $0.5$. Both outputs are
finite, and both fit the requested output shape. A singleton prompt test
or a square-mask test cannot distinguish this error.

\EngineFigure{foundation-attention-offset-mask}{Logical positions determine
the permitted entries in this illustrative two-query chunk. Numbered cells
are allowed; crosses are forbidden. The wrong panel treats local row zero
as logical position zero. Colour is redundant with cell labels.}
{fig:foundation-attention-offset-mask}

Let $S_{\rm prefix}$ be the number of cached positions before a new chunk, $C$ its
number of tokens, $r\in\{0,\ldots,C-1\}$ a local query row, and
$j\in\{0,\ldots,S_{\rm prefix}+C-1\}$ a key position. With no padding or eviction,
the permitted set is
\[
 {\cal J}_r=\{j:0\le j\le S_{\rm prefix}+r\}.
\]
For a single new token, $C=1$, every key in the $S_{\rm prefix}+1$-entry history is
permitted. For an uncached prompt, $S_{\rm prefix}=0$, the familiar square triangle
returns. One expression covers both cases without assigning causal
meaning to the array's local row number.

Our decoder in the next chapter appends one token at a time. The separate
\texttt{chunk\_attention} lesson teaches the offset rule; it is not a new
chunked decoder implementation. Keys and values supplied to it include
both the prefix and the entire new chunk. The function selects the proper
prefix for each query, so an already-created future K/V cannot enter its
denominator:
\begin{minipage}{\linewidth}
\CodeLines{../code/foundations/decoder_trace.py}{8}{20}
\end{minipage}
Run \texttt{python3 code/foundations/decoder\_trace.py}. The first two
output records state positions three and four, their permitted key lists
and outputs $[1.5]$ and $[2.0]$. Change the prefix to zero without removing
the old keys and the interface rejects the inconsistent history rather
than silently choosing a different mask.

\section{Keep sequence, head and coordinate axes separate}
There are three different kinds of index in an attention calculation.
A position index chooses a token within one request. A head index chooses
a projected representation. A coordinate index chooses one feature
inside that head. Mixing these axes can preserve the total element count
while changing every comparison.

For a rectangular batch with $B$ requests and $S$ positions per request,
one possible logical layout is
\[
 \mathbf Q\in\mathbb R^{B\times H_q\times S\times d_h},\quad
 \mathbf K,\mathbf V\in
 \mathbb R^{B\times H_{kv}\times S\times d_h}.
\]
This is a declaration of axes, not a required physical stride order.
Our list implementation keeps request state separately and nests each
cached row by KV head, then coordinate. An optimized tensor may store a
different order. Read the accessor, not the total size, to establish what
an address means.

For contiguous groups, let $G=H_q/H_{kv}$ be a positive integer.
Query head $h$ reads KV head $\lfloor h/G\rfloor$. At $H_q=4$,
$H_{kv}=2$, heads zero and one read KV head zero; heads two and three
read KV head one. At $H_q=H_{kv}$ every query head has its own KV head.
At $H_{kv}=1$ all query heads share one history. In all cases the score
and normalization are computed separately for each query head.

Do not average the queries first. In general, softmax of the average
score row is not the average of the two softmax rows. The nonlinear
normalization makes those different computations. Sharing K/V storage
reduces the number of stored vectors; it does not license removing query
heads or averaging their outputs before the output projection.

Requests also need separate permitted sets. Packing tokens from several
requests into one array does not make them one causal sequence.
A batch dimension or an explicit sequence-boundary map must prevent a
query for request B from reading request A's values. Padding is storage,
not a real token: it must not acquire probability mass. The tiny routine
uses neither packed requests nor padding, so its ordinary prefix slice
cannot be reused unchanged for those layouts.

\section{Count a row before counting a kernel}
At one query head, let $N$ be the number of permitted keys.
Computing scores performs $N$ dot products of width $d_h$.
Forming the value mixture performs $d_h$ reductions over $N$ terms
when value width equals head width. Under the conventional two-FLOPs
per multiply--add count, these two products total $4Nd_h$ FLOPs.
Score scaling, exponentials, normalization and mask handling are outside
that count. It describes matrix arithmetic, not CPU instructions.

For $S=4$, causal row lengths are $1,2,3,4$: ten query--key pairs
per head. Their product work is $40d_h$ conventional FLOPs.
A dense product that computes all sixteen scores and only then masks
them performs more work than a triangular evaluation. A diagram of the
mathematical permitted set is not evidence that a kernel skips the
forbidden products.

\EngineFigure{foundation-attention-row-storage}{The executable teaching
row reads a permitted K/V history and produces one output. The scratch
score/probability rows are temporary; the history survives for later
queries. Counts assume equal key/value width and packed elements of
$b$ bytes, not Python object allocation.}
{fig:foundation-attention-row-storage}

The row reference keeps $N$ scores, then probabilities, and returns
$d_h$ output coordinates. Their lifetime differs from the $2Nd_h$
K/V elements. Scratch belongs to a calculation; history belongs to the
evaluated prefix. A fused kernel may stream through blocks and retain
less scratch. That is the bridge to \Chref{flashattention}: avoiding a
large stored intermediate does not abolish the permitted comparisons.

Reading the same KV head for multiple query heads creates an opportunity
for reuse, not a guaranteed traffic reduction. Whether hardware actually
reuses those bytes depends on the loop order, tiling and memory hierarchy.
Here we count logical operands. We do not infer DRAM bytes from Python
list traversal or advertise a speedup from the smaller head count.

\section{Use answers that do not depend on a second copy}
An independent oracle is a calculation whose failure modes differ from
the program under test. The new tests evaluate several asymmetric score
rows using 70-digit Decimal arithmetic: dot products, scaling,
exponentials and value reductions are written without calling the model's
operators. Python float results agree within the stated $10^{-12}$
absolute tolerance on these small fixtures. This is finite-case evidence,
not a proof for all magnitudes or every low-precision kernel.

Structural properties supply other checks. Probabilities are nonnegative
and sum to one within numerical tolerance. Each output coordinate lies
between the minimum and maximum permitted value coordinates: it is a
convex combination. Add the same vector $\mathbf c$ to every permitted
value, leaving keys unchanged, and the output must increase by
$\mathbf c$, because the probabilities sum to one.

These properties apply before the attention output projection and residual
addition. A learned output projection can produce values outside those
coordinate bounds. Checking the convex bound on an entire decoder layer
would reject a correct implementation. Likewise, changing values alone
is different from changing an input token, which changes Q and K too.
Tests must identify the operator boundary they claim to check.

Stable softmax does not repair an overflowing dot product upstream.
The model rejects nonfinite score rows, but it is not a general tensor
validator for every intermediate. Use finite values of controlled
magnitude in these fixtures. Chapter~\ref{app:numerics} develops the
rounding and range questions needed for actual low-precision execution.

\Needspace{10\baselineskip}
\section{Build the row, then build the mask}
The executable lesson uses ordinary Python lists:
\begin{verbatim}
python3 code/foundations/lesson.py attention
python3 -m unittest discover -s code/foundations -p 'test_*.py'
\end{verbatim}
It prints probabilities near $[0.25,0.75]$ and output near $[0.5,3]$.
The small floating-point difference from an exact fraction is expected.
Read \texttt{attention\_row} in \texttt{code/foundations/model.py}.
It checks widths, calculates scores, applies stable softmax and accumulates
one output coordinate at a time. Then read \texttt{causal\_attention}.
Its prefix slice is the entire causal policy.

Before running a change, predict which output row should move. Replacing
the second value changes the second output, but must not change the first.
Giving the future key an enormous score is an even stronger test: an
unmasked implementation will attend almost entirely to it. The committed
test includes this deliberately tempting future and checks that the first
row remains exactly its own value.

Uniform probabilities provide another oracle. If all permitted scores
are equal, attention returns the arithmetic mean of their values. A
single permitted value must pass through unchanged. These are independent
facts about the operation, not a comparison against another copy of the
same implementation.

\section{What caching can change, and what it cannot}
When one more token arrives, old positioned keys and values need not be
recomputed for an unchanged causal prefix and fixed model configuration.
The new query still compares against all permitted cached keys. Caching
removes repeated projection and earlier-layer work; it does not remove
the weighted sum over the history.

Our complete decoder will keep a separate history for every layer.
The current token's key and value must join that layer's history before
its own attention is evaluated, because ordinary causal attention includes
the current position. A sampled output token is different: it has not
been evaluated yet. Sampling an ID does not create its K/V vectors.

For this chapter, storing an $S\times S$ score matrix is acceptable as a
description of the operation. It becomes expensive at long context.
\Chref{flashattention} later changes the evaluation order to avoid writing
that entire matrix. The mathematical result remains the reference we have
just built. Understanding the row first makes that optimization a question
about storage and reduction, rather than a new kind of attention.

\section{What we have built}
We can now turn positioned Q, K and V into a contextual vector, explain
every weight in a small example, and prove that a future token cannot
affect an earlier row. Attention mixes information across positions.
The next chapter supplies a different operation: transforming the features
of each position without mixing positions at all.
\begin{ConstraintLedger}
Compute & spends & Permitted pairs still require score and value products; causal layout alone does not establish skipped kernel work.\\
Memory capacity & spends & K/V persists; score and probability rows are temporary reference scratch.\\
Memory bandwidth & moves & GQA permits K/V reuse, but loop and hardware determine realized traffic.\\
Correctness & risks & Offset masks, sequence boundaries and head mapping must preserve the same permitted set.\\
\end{ConstraintLedger}
\Needspace{10\baselineskip}
\section{Worked problems}
\subsection*{1. Equal scores}
\textbf{Problem.} Two permitted values are $[2,0]$ and $[0,4]$, and their
scores are equal. Find the output.
\textbf{Solution.} Equal exponentials give weights $[1/2,1/2]$ and output
$[1,2]$. Equality of scores matters; their common finite value does not.
\Needspace{9\baselineskip}
\subsection*{2. A mask applied too late}
\textbf{Problem.} A first-position query has unmasked probabilities
$[1/4,3/4]$. Why is zeroing the second probability insufficient?
\textbf{Solution.} The retained mass is $1/4$, not one. Multiplying values
gives $[0.5,0]$ instead of $[2,0]$. Exclude the future score before softmax,
or renormalize the permitted distribution.
\subsection*{3. Test the causal claim}
\textbf{Problem.} Write a test that changes only the last token's value
and checks every earlier output row.
\textbf{Solution sketch.} Compute the original outputs, change that value
to a large vector, recompute, and compare rows before the last. None may
change. Also demonstrate that removing the prefix slice makes the test fail.
\Needspace{8\baselineskip}
\subsection*{4. Two new queries after a prefix}
\textbf{Problem.} Three positions are cached and two new query rows are
supplied with the full five-entry K/V history. Scores are equal and values
are $[0],[1],[2],[3],[4]$. Find the permitted sets and outputs.
\textbf{Solution.} The local rows have logical positions three and four.
Their permitted sets are $\{0,1,2,3\}$ and $\{0,1,2,3,4\}$, giving
$[1.5]$ and $[2]$. A local upper-left triangle gives $[0]$ and $[0.5]$,
which are shape-correct wrong answers.
\Needspace{8\baselineskip}
\subsection*{5. Shared storage, separate queries}
\textbf{Problem.} Four query heads share two KV heads in contiguous groups.
Which KV head does query head three read? How many permitted pairs do the
chapter's two-row chunk calculate across all four query heads?
\textbf{Solution.} Group size is $4/2=2$, so head three reads
$\lfloor3/2\rfloor=1$. Each query head has $4+5=9$ pairs;
the four heads calculate 36 pairs. KV sharing does not divide that count.
\Needspace{8\baselineskip}
\subsection*{6. Translate every value}
\textbf{Problem.} Add $[10,-4]$ to every permitted value in a row.
Keys and query are unchanged. What happens to the output?
\textbf{Solution.} Probabilities stay unchanged. Their sum is one, so
the output gains exactly $[10,-4]$ in real arithmetic. Floating-point
tests use a tolerance. This property applies before the output projection.

